% ============================================================
%  Chapter 11 — Dimension
% ============================================================
\section{Dimension}

The \emph{dimension} of a vector space is the number of vectors in any basis.
For this to be a well-defined concept we must prove that all bases have the
same size — a consequence of the Fundamental Lemma.  Dimension is the single
most important numerical invariant of a finite-dimensional vector space.

% ---- Definition --------------------------------------------
\begin{defbox}[Definition --- Dimension of a Vector Space]
\label{def:dimension}
Let $V$ be a finite-dimensional vector space over $\F$ with basis
$B=\{v_1,\ldots,v_n\}$.  The \textbf{dimension} of $V$ over $\F$ is
\[
    \dime_\F V = n.
\]
The trivial space $\{\mathbf{0}\}$ has dimension $0$.  If $V$ has no finite
basis it is called \textbf{infinite-dimensional}.
\end{defbox}

% ---- Fundamental Lemma -------------------------------------
\begin{thmbox}[Theorem --- Fundamental Lemma of Linear Algebra]
\label{thm:fundamental-lemma}
If $\{v_1,\ldots,v_n\}$ is a linearly independent set in $V$, then every
subset of $\spn\{v_1,\ldots,v_n\}$ with more than $n$ vectors is linearly
dependent.
\end{thmbox}

\begin{proofbox}[Proof --- Fundamental Lemma of Linear Algebra]
Suppose $w_1,\ldots,w_m\in\spn\{v_1,\ldots,v_n\}$ with $m>n$.  Write
$w_j=\sum_{i=1}^n a_{ij}v_i$ for each $j$.  The system
$\sum_{j=1}^m\alpha_j w_j=\mathbf{0}$ becomes
$\sum_{i=1}^n\!\bigl(\sum_{j=1}^m a_{ij}\alpha_j\bigr)v_i=\mathbf{0}$.
Since $\{v_i\}$ is l.i., this reduces to the homogeneous system
$(a_{ij})\alpha=\mathbf{0}$ with $n$ equations in $m>n$ unknowns, which
always has a non-trivial solution.
\hfill$\blacksquare$
\end{proofbox}

% ---- Dimension is well-defined -----------------------------
\begin{thmbox}[Theorem --- Dimension Is Well-Defined]
\label{thm:dimension-well-defined}
Any two bases of a finite-dimensional vector space $V$ have the same number of
elements.  Hence $\dime_\F V$ is well-defined.
\end{thmbox}

\begin{proofbox}[Proof --- Dimension Is Well-Defined]
Let $B=\{v_1,\ldots,v_n\}$ and $B'=\{w_1,\ldots,w_m\}$ both be bases.
Since $B$ is l.i.\ and $B'\subseteq\spn(B)=V$, the Fundamental Lemma gives
$m\le n$.  By symmetry (swapping $B$ and $B'$) we also get $n\le m$, so
$m=n$.
\hfill$\blacksquare$
\end{proofbox}

% ---- Standard dimensions -----------------------------------
\begin{exbox}[Example --- Dimensions of Standard Spaces]
\begin{enumerate}
    \item $\dime_\F \F^n = n$ (standard basis $e_1,\ldots,e_n$).
    \item $\dime_\F \Mmn{m}{n} = mn$ (basis $\{E^{ij}\}$).
    \item $\dime_\R \C^1 = 2$, since $\C\cong\R^2$ as real vector spaces
          (basis $\{1,i\}$ over $\R$).
    \item $\dime_\R \C^n = 2n$.
\end{enumerate}
\end{exbox}

% ---- Column space, row space, rank -------------------------
\begin{defbox}[Definition --- Column Space and Row Space]
\label{def:column-row-space}
For $A=[v_1\mid\cdots\mid v_n]\in\Mmn{m}{n}$ with columns $v_i\in\F^m$ and
rows $w_j\in\F^n$:
\begin{align*}
    \text{Column space of }A &= \spn\{v_1,\ldots,v_n\} \subseteq \F^m. \\
    \text{Row space of }A    &= \spn\{w_1,\ldots,w_m\} \subseteq \F^n.
\end{align*}
\end{defbox}

\begin{defbox}[Definition --- Rank of a Matrix]
\label{def:rank-matrix}
The \textbf{rank} of $A\in\Mmn{m}{n}$ is
\[
    \rk(A) = \dime_\F(\text{column space of }A).
\]
\end{defbox}

% ---- Computing rank ----------------------------------------
\begin{tipbox}[Remark --- Computing Rank]
\label{rem:rank-rref}
$\rk(A)$ equals the number of nonzero rows in the RREF of $A$, which also
equals the number of pivot columns.  In particular:
\begin{itemize}
    \item Row-equivalent matrices have the same rank.
    \item $\rk(A) = \rk(R)$ where $R$ is the RREF of $A$.
\end{itemize}
\end{tipbox}

% ---- Subspace dimension ------------------------------------
\begin{thmbox}[Theorem --- Dimension of a Subspace]
\label{thm:subspace-dim}
If $W$ is a subspace of a finite-dimensional vector space $V$, then $W$ is
also finite-dimensional and
\[
    \dime_\F W \;\le\; \dime_\F V.
\]
Moreover $\dime_\F W = \dime_\F V$ if and only if $W=V$.
\end{thmbox}

\begin{proofbox}[Proof --- Dimension of a Subspace]
Start with $\mathbf{0}\in W$ (or the empty set if $W=\{\mathbf{0}\}$).
Repeatedly pick vectors in $W$ not in the span of those chosen so far.  By the
Fundamental Lemma, this process terminates after at most $n=\dime V$ steps,
yielding a finite l.i.\ spanning set for $W$, so $W$ is finite-dimensional and
$\dime W\le n$.  If $\dime W=n$ then $W$ has a basis of $n$ l.i.\ vectors,
which also spans $V$ by Theorem~\ref{thm:basis-char}, so $W=V$.
\hfill$\blacksquare$
\end{proofbox}
